% ============================================================
% 油气田生产网 IO 服务器 — 系统本体
% V1.0 · 2026-07-11
% ============================================================
\documentclass[12pt,a4paper]{ctexart}
\usepackage[top=2cm,bottom=2cm,left=2cm,right=2cm]{geometry}
\usepackage{graphicx,float,longtable,booktabs,array,enumitem}
\usepackage{fancyhdr,hyperref}
\usepackage{xcolor,tikz}
\usetikzlibrary{positioning,arrows,shapes}
\setmainfont{Microsoft YaHei}
\setCJKmainfont{Microsoft YaHei}
\setmonofont{Consolas}
\pagestyle{fancy}
\fancyhf{}
\fancyhead[L]{IO 服务器本体 — 油气田生产网}
\fancyhead[R]{V1.0 · 2026-07-11}
\fancyfoot[C]{\thepage}
\renewcommand{\headrulewidth}{0.4pt}
\hypersetup{colorlinks=true,linkcolor=blue,pdftitle={IO服务器本体}}

\begin{document}

\begin{titlepage}
\centering
\vspace*{3cm}
{\Huge\bfseries 油气田生产网 IO 服务器}\\[0.5cm]
{\LARGE 系统本体}\\[2cm]
{\large 迪格（杭州）物联科技有限公司}\\[0.3cm]
{\large 2026 年 7 月 11 日}\\[1cm]
\begin{tabular}{rl}
\textbf{版本} & V1.0 \\
\textbf{状态} & 调研完成 \\
\textbf{密级} & 内部 \\
\end{tabular}
\end{titlepage}

\tableofcontents
\newpage

% ============================================================
\section{硬件配置}
% ============================================================

\begin{table}[H]\centering
\caption{IO 服务器硬件}
\begin{tabular}{p{3cm}p{5cm}p{5cm}}
\toprule
\textbf{项目} & \textbf{IO 服务器 (131)} & \textbf{pSpace 协议端口 (130)} \\
\midrule
主机名 & WIN-NJ7VP96R2NM & --- \\
操作系统 & Windows 10.0.14393 & --- \\
CPU & Intel Xeon (推测) & --- \\
网卡 & Intel X520 10GbE ×4 + X722 1GbE ×4 & --- \\
IP & 11.66.12.131 & 11.66.12.130 \\
内存 & 32GB+ (pSpace 进程占用 $\sim$1.2GB) & --- \\
磁盘 & E: 1.1TB 空闲 & --- \\
\bottomrule
\end{tabular}
\end{table}

% ============================================================
\section{软件栈}
% ============================================================

\subsection{pSpace 平台}

\begin{table}[H]\centering
\caption{pSpace 组件}
\begin{tabular}{p{3cm}p{1.5cm}p{2cm}p{1.5cm}p{4cm}}
\toprule
\textbf{进程} & \textbf{实例数} & \textbf{单实例内存} & \textbf{版本} & \textbf{功能} \\
\midrule
IoProject.exe & 1 & 306 MB & --- & IO 项目主引擎 \\
IOMan.exe & 36 & $\sim$16 MB & 7.0.3.2 & IO 设备管理器 \\
IoMonitor.exe & 1 & 61 MB & 7.0.1.2 & 监控界面 \\
CommBridge.exe & 1 & 32 MB & 6.3.0.1 & 通讯网桥 \\
IoCommit.exe & 7 & $\sim$15 MB & --- & 数据提交 \\
\bottomrule
\end{tabular}
\end{table}

\subsection{依赖环境}

\begin{itemize}[leftmargin=*]
  \item Oracle Client 11.2.0 (E:\textbackslash app\textbackslash Administrator\textbackslash product\textbackslash 11.2.0\textbackslash client\_1)
  \item OPC Core Components (32位, SysWOW64)
  \item .NET Framework 4.0
  \item 力控 pSpace SDK (psAPISDK.dll, 3525 导出函数)
\end{itemize}

\subsection{安装路径}

\begin{verbatim}
  E:\IO ServerOnLine\
  ├── IoProject.exe      (主引擎)
  ├── IOMan.exe ×36      (设备管理器)
  ├── IoMonitor.exe       (监控界面)
  ├── CommBridge.exe      (通讯网桥)
  ├── IoCommit.exe ×7     (数据提交)
  ├── psAPISDK.dll        (力控 API, 3.5MB, 32位)
  ├── Device.ini          (12种保护装置配置)
  ├── IoChannelCfg.ini    (通道配置)
  ├── IoMonitor.ini       (提交参数)
  ├── SqlFilSet.ini       (Oracle 连接串)
  ├── Event.txt           (历史事件日志)
  └── 11.66.11.131\       (配置备份)
       └── HisData\        (历史数据目录)
\end{verbatim}

% ============================================================
\section{网络拓扑}
% ============================================================

\begin{figure}[H]\centering
\begin{tikzpicture}[
  box/.style={draw,rounded corners,minimum width=3cm,minimum height=1.2cm,align=center,font=\small},
  arrow/.style={->,>=stealth,thick,font=\footnotesize}
]
  \node[box,fill=blue!15] (io131) at (0,0) {\textbf{11.66.12.131}\\IO 服务器\\pSpace 平台};
  \node[box,fill=green!12] (pspace) at (-5,2) {\textbf{11.66.12.130}\\pSpace A11\\:8889};
  \node[box,fill=gray!10] (oracle) at (5,2) {\textbf{11.66.12.129}\\Oracle 功图库\\:1521};
  \node[box,fill=orange!10] (dcs) at (-5,-2) {DCS RTU/PLC\\172.x (5套)\\OPC DA (RSLinx)};
  \node[box,fill=orange!10] (rtu) at (0,-2.5) {Modbus RTU\\11.248-250.x\\206台};
  \node[box,fill=yellow!12] (gprs) at (5,-2) {无线终端\\GPRS/CDMA\\47台};
  \node[box,fill=green!8] (dev) at (0,-4.5) {\textbf{11.66.191.155}\\开发机\\dgiot\_lite};

  \draw[arrow] (pspace) -- node[right]{A11 5a5a} (io131);
  \draw[arrow] (io131) -- node[above]{TNS} (oracle);
  \draw[arrow] (io131) -- node[left]{DCE/RPC OPC DA} (dcs);
  \draw[arrow] (io131) -- node[right]{Modbus TCP :53001} (rtu);
  \draw[arrow] (io131) -- node[right]{UDP} (gprs);
  \draw[arrow,red,thick] (dev) -- node[right]{WinRM :5985} (io131);
\end{tikzpicture}
\caption{IO 服务器网络拓扑}
\end{figure}

% ============================================================
\section{核心 IP / 端口 / 凭据}
% ============================================================

\begin{table}[H]\centering
\caption{核心节点}
\begin{tabular}{p{3cm}p{3.5cm}p{2cm}p{4.5cm}}
\toprule
\textbf{IP} & \textbf{角色} & \textbf{端口} & \textbf{说明} \\
\midrule
11.66.12.131 & IO 服务器 & 5985 & WinRM (administrator/GKYWB-5991792\ldots) \\
 & & 53001 & Modbus TCP 主站 \\
 & & 2500 & RDP 远程桌面 \\
 & & 54531-8 & IOMan 内部端口 \\
11.66.12.130 & pSpace A11 & 8889 & A11 IO 协议 \\
11.66.12.129 & Oracle & 1521 & 功图库 (DQYTPROD@orcl) \\
11.66.113.78 & DMZ 中转 & 1883/59660 & MQTT/SSH \\
11.66.191.155 & 开发机 & -- & dgiot\_lite \\
\bottomrule
\end{tabular}
\end{table}

\begin{table}[H]\centering
\caption{DCS 系统 (OPC DA)}
\begin{tabular}{p{3cm}p{3cm}p{2.5cm}p{3.5cm}}
\toprule
\textbf{IP} & \textbf{子网} & \textbf{设备数} & \textbf{OPC Server} \\
\midrule
172.23.9.3 & DCS-A & 30+ & RSLinx Classic \\
172.23.18.194 & DCS-B & 30+ & RSLinx Classic \\
172.26.6.3 & DCS-C & 60+ & RSLinx Classic \\
172.21.14.192 & DCS-D & 40+ & RSLinx Classic \\
172.28.5.200 & DCS-E & 40+ & RSLinx+WinCC \\
\bottomrule
\end{tabular}
\end{table}

\begin{table}[H]\centering
\caption{RTU 设备 (Modbus TCP)}
\begin{tabular}{p{3cm}p{2cm}p{2.5cm}p{4cm}}
\toprule
\textbf{网段} & \textbf{子网数} & \textbf{设备数} & \textbf{从站} \\
\midrule
11.248.x.x & 10 & 139 & slave=1,2 \\
11.249.x.x & 26 & 89 & slave=1,2 \\
11.250.x.x & 22 & 72 & slave=1 \\
\bottomrule
\end{tabular}
\end{table}

% ============================================================
\section{数据流}
% ============================================================

\begin{verbatim}
  RTU/PLC (Modbus TCP :502)
      │
      ▼
  IOMan ×36 ──→ IoProject ──→ IoCommit ×7 ──→ Oracle (129:1521)
      │              │                │
      │ pSpace API   │ 事件流         │ 功图数据
      │              ▼                ▼
      │         IoMonitor         DQYTPROD@orcl
      │         (GUI 显示)
      │
  CommBridge ←── GPRS/CDMA 无线终端 (47台)
      │
      ▼
  数据转发到 DMZ
\end{verbatim}

% ============================================================
\section{协议栈}
% ============================================================

\begin{table}[H]\centering
\caption{协议一览}
\begin{tabular}{p{2.5cm}p{5cm}p{1.5cm}p{3.5cm}}
\toprule
\textbf{协议} & \textbf{通信对} & \textbf{帧数} & \textbf{状态} \\
\midrule
A11 (5a5a) & 130:8889 $\leftrightarrow$ 131 & 93.9K & 已逆向，含设备路径 \\
Modbus TCP & 131:53001 $\rightarrow$ RTU & 21.3K & 已解析，206台寄存器值 \\
OPC DA (DCE/RPC) & 172.x $\leftrightarrow$ 131 & 72.9K & RSLinx/WinCC，DCOM 加密 \\
Oracle TNS & 131:53268 $\rightarrow$ 129:1521 & 8.7K & 功图数据查询 \\
\bottomrule
\end{tabular}
\end{table}

\subsection{A11 帧结构}

\begin{verbatim}
  Offset  Size   Field         Value
  0       2B     Magic         5a 5a
  2       2B     Frame Length  LE, 不含 2B 头
  4       4B     Flags         控制标志
  8       2B     Message Type  0xF062=设备列表, 0x87B2=心跳
  10      N      Payload       ASCII 设备路径 + 浮点值
\end{verbatim}

\subsection{Modbus TCP 寄存器}

\begin{table}[H]\centering
\begin{tabular}{p{2cm}p{3cm}p{4cm}p{4cm}}
\toprule
\textbf{地址} & \textbf{参数} & \textbf{量程公式} & \textbf{示例值} \\
\midrule
299-302 & 模拟量 & Y $\times$ 170 / 8192 (A) & 15930 $\rightarrow$ 330A \\
980 & 状态量 & 原始值 & 60 \\
7443 & 批量模拟量 & Y $\times$ 170 / 8192 (V) & 4960 $\rightarrow$ 103V \\
\bottomrule
\end{tabular}
\end{table}

% ============================================================
\section{事件日志样本（IoMonitor 实时数据）}
% ============================================================

\begin{verbatim}
  02110120089 B1V25VE33: 最大载荷点=86.070000
  02110150041 B1V51VSFK01: 最大下行电流点=15.280000
  02110150041 B1V51VSFK01: 最小载荷点=33.530000
  02110150041 B1V51VSFK01: 最大上行电流点=10.690000
  02110150041 B1V51VSFK01: 最大载荷点=41.310000
  02110150041 B1V51VSFK01: =====插入数据库成功=====
  21001086003: 打开CommBridge通道成功，编号409
  02204060071: 打开CommBridge通道失败，编号236
\end{verbatim}

% ============================================================
\section{数据采集方案}
% ============================================================

\begin{table}[H]\centering
\caption{采集路径评估}
\begin{tabular}{p{2.5cm}p{2cm}p{2.5cm}p{5cm}}
\toprule
\textbf{数据源} & \textbf{设备数} & \textbf{可行性} & \textbf{当前状态} \\
\midrule
Modbus TCP & 206 & ✅ 可实时采 & 已模拟 206 台，明文数据 \\
A11 & --- & ✅ 可解析 & pcap 中含设备路径 \\
OPC DA & 200+ & ⚠️ DCOM & 需配置 DCOMCNFG \\
Oracle 功图 & 全量 & ⚠️ 密码 & DQYTPROD@orcl \\
IoMonitor GUI & 全量 & ✅ 可读 & UI Automation 直接读事件文字 \\
\bottomrule
\end{tabular}
\end{table}

% ============================================================
\section{开发机环境}
% ============================================================

\begin{table}[H]\centering
\begin{tabular}{p{2cm}p{4cm}p{6cm}}
\toprule
\textbf{端口} & \textbf{服务} & \textbf{说明} \\
\midrule
:5173 & Vue 仪表盘 & http://localhost:5173/capture \\
:8765 & capture\_server & 多协议解析 A11/Modbus/OPC-DA/IEC104 \\
:53001 & modbus\_sim & 206 台 RTU Modbus TCP 模拟器 \\
\bottomrule
\end{tabular}
\end{table}

\begin{table}[H]\centering
\caption{开发机网络}
\begin{tabular}{p{3cm}p{4cm}p{2.5cm}p{3cm}}
\toprule
\textbf{接口} & \textbf{IP} & \textbf{网关} & \textbf{用途} \\
\midrule
WLAN & 192.168.3.111 & 192.168.3.1 & 互联网 \\
以太网 & 11.66.191.155 & 11.66.191.251 & 企业内网 \\
\bottomrule
\end{tabular}
\end{table}

\subsection{路由}

\begin{verbatim}
  route add 11.66.0.0 mask 255.255.0.0 11.66.191.251 metric 5 -p
\end{verbatim}

\vspace{1cm}
\begin{center}
\rule{0.5\textwidth}{0.4pt}\\[6pt]
{\small 迪格（杭州）物联科技有限公司 · 2026 年 7 月 · V1.0}
\end{center}

\end{document}
